% lit_membership.tex -- Literature 2.2: changes of team membership, with the numeric table T3 (drafted 12 Oct 2026
% from the thesis plan of 11 Oct). The facts on Quan et al. and Zeng et al. were verified from the full texts on
% 11 Oct (hand-back docs/thesis_work_2026-10-11.md); the cells in brackets wait for a full-text read (plan C2-A:
% Bacheti and Delbene were read from abstracts only) or for the October flights. Entries the plan lists for this
% section that are not yet in the bibliography (Oishi 2023, Gao 2024, Delbene 2024/2025, Bacheti 2026, Fink 2011,
% Tognon 2018, Xu 2025, Farina 2020) wait for the approval batch.
% Papers found with Undermind and approved by Wesley, 10 Oct (DOIs, abstracts and sources in
% docs/thesis_drafts/refs_2026-10-10.md; alternates listed there).

\subsection{Changes of team membership}
\label{sec:lit_team}

In this thesis a reconfiguration is a change of team membership: a vehicle leaving or joining while the payload
is carried. The word is also used for a fixed team changing its shape, as when a team contracts to pass a doorway
\cite{jackson2020scalable, tensionbased}; that is a formation change, and the team's force allocation keeps its
dimension. A membership change does not. Most of the work of Section~\ref{sec:lit_models} assumes a fixed team
\cite{efficient, wahbaKinodynamicMotionPlanning2024, quan2025swarm}; the few works that change it are set out in
Table~\ref{tab:membership}, with the detection of the change kept apart from the recovery from it, because the
two are often reported as one time.

\paragraph{A vehicle leaving, on hardware.} Quan et al.\ \cite{quan2025swarm} remove a vehicle from a carried
load in flight: four robots carry a 4\,kg load, one releases its cable on command, and the other three
reorganise to a near-equilateral layout and rebalance their tensions within 13\,s, decentralised and without
replanning. They report the cable tensions and one robot's velocity through the event, not the load's attitude.
Zeng et al.\ \cite{zeng2025marl} switch one of three vehicles off in flight and leave it hanging from a 1.4\,kg
load; the other two keep five degrees of freedom of the load under a learned decentralised policy that observes
only the load's pose, and the load's position and attitude are plotted through the loss with no transient
number. Bacheti et al.\ \cite{tensionbased} handle a removal by repositioning the remaining vehicles before the
release so that the transient is small, on hardware with a point-mass payload; [whether the agent count changes
in their experiments is to be verified from the full text]. So on hardware a vehicle has been released from a
carried load and a vehicle has been lost from one, and in neither case is the load's attitude through the event
given as a number.

\paragraph{A link that fails, in simulation.} A cable that breaks rather than releases is treated as a fault.
Liang et al.\ \cite{liang2021fault} estimate the unmeasured cable tensions with finite-time disturbance
observers and reassign the load among the remaining quadrotors after a cable breaks; Saiella et al.\
\cite{saiella2021fault} let a supervisor reconfigure the network topology and the control objectives of a
passivity-based formation when a vehicle faults; Hajieghrary and Schmitt \cite{hajieghrary2026passive} route
each vehicle's measured cable tension straight into its thrust command, so that five vehicles recover from one or
two severed cables within a payload-pendulum period, with a hybrid input-to-state stability certificate; and
Farras et al.\ \cite{farrasFaultTolerantControlMultiQuadcopter2024} keep a team with a suspended payload flying
through a fault under wind. None of the four flies, where a time is given it is the recovery and not the
detection of the loss, and none adds a vehicle.

\paragraph{A vehicle joining.} No work attaches a vehicle to a payload that a team is already carrying, in
simulation or on hardware. Quan et al.\ describe a joining robot only in prose, as one already connected to the
payload and its neighbours, and give it no mechanism to learn its share of the load or its attachment point
\cite{quan2025swarm}. The event-triggered architecture of Tasooji et al.\
\cite{tasoojiCooperativeControlMultiQuadrotors} is argued to scale to more vehicles, but a varying team size is
not analysed and its join protocol has neither an experiment nor an explicit force reallocation. The closest
hardware precedent is a payload handed from one multirotor to another: Shankar et al.\ \cite{shankar2020inflight}
grasp a bar held aloft by a second vehicle, confirm the grasp with a magnetic sensor before the first releases,
and reschedule the gains for the new mass; the team does not grow, and the payload is not shared.

\paragraph{Prepared changes: plug-and-play control.} Plug-and-play model predictive control lets subsystems join
or leave a running network without losing stability or constraint satisfaction, provided each change is
prepared before it happens. Riverso et al.\ \cite{riverso2013plug} redesign offline only the controllers of the
subsystems a newcomer couples to and refuse a change the redesign cannot certify, Zeilinger et al.\
\cite{zeilinger2013plug} and Aboudonia et al.\ \cite{aboudonia2022reconfigurable} add a transition phase that
prepares the network before the switch, and Saccani et al.\ \cite{saccani2023model} remove the request by having
each robot also plan a safe trajectory that holds if others join within the horizon [Farina, Zhang and Scattolini
2020, pending the approval batch]. All four are shown in simulation, on linear coupled subsystems such as power
networks and coupled masses, or on ground robots coupled only through collision constraints. A suspended payload
couples every vehicle through the load, and a magnet that lets go or a contact that catches happens at an instant
the planner does not choose; this thesis therefore switches a centralised nonlinear problem to a pre-built one of
the new size at the event, warm-started from the measured state, and shows the first solve after the switch in
flight rather than certifying it (Section~\ref{sec:reconfig}).

\paragraph{What statics decides.} Static equilibrium decides which layouts can hold a suspended payload at all.
Michael et al.\ \cite{michael2011cooperative} show that three robots are the fewest that hold a payload over a
wide set of poses, require the attachment points to span its centre of mass, and find several equilibria of which
only some are stable; Jiang and Kumar \cite{jiang2013inverse} show that the robot positions for a given pose are
not unique for three or more robots and become finitely many for three once the tensions are fixed. Gabellieri et
al.\ \cite{gabellieri2023equilibria} add that the internal force of non-vertical cables makes the load's attitude
equilibrium stable and less sensitive to model error, for two robots and a beam on hardware, and Liu et al.\
\cite{liu2025wrench} compute from the same balance how much wrench a platform on rigid links can still receive
within the vehicles' thrust limits [Fink 2011 and Tognon 2018, pending the approval batch, for the planar case and
its qualifier]. All four analyse a fixed team: none asks what the equilibrium becomes when one carrier leaves the
layout in flight. Section~\ref{sec:reconfig} takes the condition that the centre of mass stay inside the
survivors' polygon and turns it into a margin that says which carrier a layout may lose and by how much.

\begin{table}[htbp]
\centering
\footnotesize
\setlength{\tabcolsep}{3pt}
\caption[Changes of team membership under a carried payload]{Changes of team membership under a carried
payload, as each work reports them. Announced: the controller is told of the change before it happens. Leaver:
whether the vehicle leaves the load or stays attached to it. Detection is the time from the event to the
controller acting on it; recovery is the time to a settled state. ``--'': not reported. Cells in brackets are to be
filled from the full texts or from the October flights.}
\label{tab:membership}
\begin{tabularx}{\textwidth}{@{}
  >{\raggedright\arraybackslash\hsize=0.86\hsize}X
  >{\raggedright\arraybackslash\hsize=1.04\hsize}X
  >{\raggedright\arraybackslash\hsize=0.78\hsize}X
  >{\raggedright\arraybackslash\hsize=0.90\hsize}X
  >{\raggedright\arraybackslash\hsize=1.04\hsize}X
  >{\raggedright\arraybackslash\hsize=1.14\hsize}X
  >{\raggedright\arraybackslash\hsize=1.24\hsize}X@{}}
\toprule
Work & Event; mechanism & Announced & Leaver & Plant; team; payload & Attitude or sag through the event &
Detection; recovery \\
\midrule
Quan\newline \cite{quan2025swarm}, 2025 & release on command; cable released at the load & yes & removed, flies off with
its cable & hardware, outdoor; 4 to 3; 4\,kg & -- (tensions and one robot's velocity plotted) & --; tensions
balanced after 13\,s \\
\addlinespace
Zeng\newline \cite{zeng2025marl}, 2025 & motors switched off & no & left hanging from the load & hardware; 3 to 2;
1.4\,kg & plotted, no number & none explicit (learned policy); -- \\
\addlinespace
Bacheti\newline \cite{tensionbased}, 2023 & planned release; the team repositions first & yes & [removed; count to
verify] & hardware; [3 to 2 to 3]; point mass & -- & --; -- \\
\addlinespace
Liang\newline \cite{liang2021fault}, 2021 & cable breaks & no & removed & simulation; [team]; [mass] & [from the paper] &
disturbance observers; [from the paper] \\
\addlinespace
Saiella\newline \cite{saiella2021fault}, 2021 & vehicle fault; the supervisor reconfigures & no & [from the paper] &
simulation; [team] & -- & supervisor; -- \\
\addlinespace
Hajieghrary\newline \cite{hajieghrary2026passive}, 2026 (preprint) & one or two cables severed & no & removed &
simulation; 5 to 4 and 5 to 3 & [from the preprint] & none (passive feed-forward); within a pendulum period
[0.32\,s] \\
\addlinespace
Farras\newline \cite{farrasFaultTolerantControlMultiQuadcopter2024}, 2024 & [fault, from the paper] & no & [from the
paper] & simulation; [team] & [from the paper] & [from the paper] \\
\addlinespace
Shankar\newline \cite{shankar2020inflight}, 2020 & a bar handed from one multirotor to another & yes & the giver releases
after the grasp is confirmed & hardware; 1 to 1 & -- & magnetic sensor at the grasp; gains rescheduled \\
\addlinespace
ModQuad\newline \cite{saldana2018modquad}, TRADY \cite{sugihara2023trady}, LEGION \cite{sugihara2026legion}, Miyazaki \cite{miyazaki2018airborne} & vehicles docking to one another; no payload & yes & n/a & hardware; 1 to 1 per dock &
TRADY: altitude jump 0.39\,m, 0.065\,m with the thrust blend & motion capture, contact force or a latch; -- \\
\midrule
This work & magnet switched off, commanded or not (leave); contact and a proven lift (join) & both & removed; the
freed UAV lands by itself & hardware and rig twin; 4 to 3 and 3 to 4; 0.86\,kg ring on rigid rods & [8 October:
tilt rise 5--17$^\circ$, absolute peak 6.5--17.8$^\circ$, dip 3.7--6.9\,cm over four releases; extended by the
14 October flights] & 0.33--0.42\,s from the rod's stretch; [time back in band, 14 October]; join: [the claim
table of 16 October] \\
\bottomrule
\end{tabularx}
\end{table}
